\begin{document}
\twocolumn[
\aistatstitle{Adaptive Interaction Graphs for Particle Simulation}
\aistatsauthor{Anonymous Authors}
\aistatsaddress{}]
\fancyhead[CE,CO]{\small\bfseries Adaptive Interaction Graphs for Particle Simulation}

\begin{abstract}
More messages do not automatically make a better particle simulator. They change the computation of a model trained on a particular interaction graph. We construct adaptive graphs that preserve native edges and append an exact budget of nearby pairs, then study how the simulator learns to use these messages and where a controller should place them. Mixed-graph training exposes the model to random expansions; cached residual scores allocate pairs with one simulator pass per forecast after initialization. Three-seed studies on Goop, WaterDrop and Sand, with paired continuations for WaterDrop, reveal different answers to these two questions. Exposure improves fixed-random rollout error in every Goop and WaterDrop seed, reducing Goop mean MSE from 0.301 to 0.129. Residual-guided placement is less reliable. On Sand, exposure lowers observed-history error under risk-guided allocation while worsening cached-risk rollout error in every seed. Prediction difficulty therefore differs from the benefit of an added interaction, and a favorable comparison on observed states need not survive autonomous feedback. These distinctions provide a practical basis for constructing and evaluating adaptive particle graphs.
\end{abstract}

\section{Introduction}
A graph-based particle simulator decides which particles exchange messages before predicting their motion\citep{sanchez2020}. Enlarging a neighborhood seems an obvious way to supply more information near difficult interactions. Yet additional messages can also perturb a prediction that was already accurate. The model has learned a computation on one distribution of graphs; changing its connectivity at prediction time changes that computation. Before asking where to add interactions, we must ask whether the simulator has learned to use them.

Residual prediction offers an appealing allocation rule: spend more computation where the model expects a large error. Its weakness is equally intuitive. A simulator may be wrong because it lacks information from another particle, because the state leaves some forcing unobserved, or because the model cannot express the required dynamics. Extra messages can address some of these errors and leave others untouched. A good prediction of difficulty need not identify a useful intervention.

We investigate this gap with constructive adaptive graphs. We retain the simulator's native graph and append a specified number of optional local pairs. Graph-exposure training lets the model encounter such expansions before evaluation. A cached residual controller then chooses pairs using scores from its preceding prediction. Figure~\ref{fig:constructive-overview} shows the graph operation and the three comparisons it enables: change training at a fixed policy, compare placement on common histories, and follow placement through autonomous feedback.

The resulting evidence contains a clear tension. Under the same random expansion policy, graph exposure improves every paired Goop and WaterDrop seed. But Goop residual allocation still loses to random on observed test histories. Sand makes the consequence visible: exposure narrows its observed risk deficit while worsening cached-risk rollouts in every seed. Our contribution is the constructive intervention and the paired study of this separation. The central question is whether a graph change improves the model's prediction, rather than whether its endpoints are difficult to predict.

% VISUAL_OVERVIEW_INSERT

\section{Constructive adaptive graphs}
\subsection{Changing messages without replacing the native graph}
The simulator receives six consecutive frames $H_t$, particle attributes and boundary information. It predicts normalized acceleration and a positive residual score for each particle; acceleration is transformed back to physical coordinates before integration. Let $E_0(x_t)$ be its native directed graph, including self-messages and the receiver cap of 128. We preserve this graph and collect optional unordered pairs in a surrounding annulus. If $P_\rho(x_t)$ contains the unordered geometric pairs inside radius $\rho$, then
\[
 C_R(x_t)=P_R(x_t)\setminus P_r(x_t),\qquad R>r.
\]
The geometric boundary convention is specified in Appendix~\ref{sec:implementation-details}.
For an integer budget $B_t\ge0$, choose $S_t\subseteq C_R(x_t)$ with $|S_t|=\min(B_t,|C_R(x_t)|)$. Each selected pair adds both orientations, $D(S_t)=\{(i,j),(j,i):\{i,j\}\in S_t\}$, so
\begin{equation}
 E_t=E_0(x_t)\cup D(S_t),\qquad |E_t|=|E_0(x_t)|+2|S_t|.
 \label{eq:budget}
\end{equation}
Our budgeted policies use $B_t=\lfloor0.25|C_R(x_t)|\rfloor$. They therefore add exactly the same number of messages when evaluated on common positions. Short-range pairs omitted by the native cap are not restored by the annulus, and additions are not recapped; final receiver degrees may exceed 128. The budget controls the intervention, not the simulator's total computational cost.

The policy determines placement within this common action space. Random25 samples uniformly; Speed25 scores displacement magnitude; Risk25 uses predicted residual scores. RMS25 uses relative velocity over native neighbors. Native base and dense expansion provide the no-addition and all-annulus controls. All six policies are used for Goop and Sand; RMS25 was not predeclared for WaterDrop. The radius ratio is 1.267. Appendix~\ref{sec:implementation-details} specifies pair ranking, tie handling and graph conventions.

\subsection{Learning to use additional messages}
Base-only training always uses the native graph. In the paired mixed-graph arm, each training example independently receives a uniformly sampled quarter-annulus expansion with probability $1/2$. The two arms share initialization or continuation parent, sampled frames, noise, architecture and optimization schedule. Thus changing the training graph is the intervention; the objective does not reward an allocation policy directly.

We use established faithful regression\citep{stirn2023} to learn residual scores without reweighting the mean predictor's squared-error gradient on matched inputs and graphs. The score $q_i$ estimates per-coordinate squared residual in normalized acceleration units, so $dq_i$ corresponds to vector squared error. It is neither an isolated epistemic-uncertainty estimate nor an estimate of the benefit of expansion. The adopted loss, stop-gradient construction, gradient handling and normalization are given in Appendix~\ref{sec:implementation-details}.

\subsection{Allocating with a cache}
The autonomous residual controller ranks each optional pair by its larger cached endpoint score. One native-base pass on the initial six observed frames initializes the cache. Each forecast then selects pairs, runs the simulator once, advances positions and caches the new scores. The budget applies from forecast 1; later forecasts require no additional scoring pass and use no future positions or labels.

The cache reuses information the simulator has already produced rather than requesting a new prediction solely for allocation. Current positions determine which pairs are eligible, while preceding scores determine their rank. This creates a temporal lag: the controller can respond to the latest available residual signal, but that signal may be stale when motion changes rapidly.

The source of the score matters. During autonomous motion it comes from the controller's own preceding selected graph. For an observed-history comparison, each policy receives the same current history and risk is instead scored on the preceding observed base history. This second setting removes differences in predicted positions and permits a direct comparison of placement. It does not reproduce the feedback encountered by the autonomous controller.

% CONFERENCE_EXPERIMENTS_MAIN_INSERT

\section{From residual risk to interaction benefit}
The experiments suggest a different target for an allocator. Residual risk estimates error under the current model and graph. For a fixed model, history $H_t$ and target $Y$, the realized benefit of adding pairs $S$ is instead
\begin{align}
 \Delta(S;H_t,Y)={}&\mathcal E(f(H_t,E_0(x_t)),Y)\nonumber\\
 &-\mathcal E(f(H_t,E_0(x_t)\cup D(S)),Y).
 \label{eq:intervention-benefit}
\end{align}
Both predictions use the same loss and target; positive benefit means expansion helped. A large residual could indicate missing interaction information, but it could also reflect ambiguity or model bias that the added pairs cannot correct. Accurate residual prediction supplies no ordering of these benefits.

Benefit belongs to the selected set, not necessarily to each edge independently. Message passing can combine or oppose information from several additions. A per-particle error score lacks both a direction for changing a prediction and an explicit measure of these interactions. The budget fixes the action size; it is the response to that action that the experiments assess.

Separate WaterDrop controls show the same mismatch under faithful and conventional heteroscedastic objectives. Risk correlates with base error but weakly with its own sparse-action benefit, and random has lower one-step error in every seed. Restoring the trained self-messages lowers error without reversing that ordering. The exact loss-change analysis separates alignment with the target from the squared prediction change (Appendix~\ref{sec:actual-action-benefit}). Conditional score-perturbation bounds require the residual score to approximate useful action values; that assumption is not supplied by residual calibration (Appendix~\ref{sec:allocation-score-conditions}).

Computation also requires a common-state comparison. Adding pairs increases the number of directed messages on that state. A smaller graph later in a rollout can result from a different predicted geometry rather than cheaper placement. Caching removes a separate subsequent scoring pass, while construction, selection and extra messages still cost work. Different geometries and shared-host execution prevent an isolated-speedup conclusion; Appendix~\ref{sec:physical-cost} gives the decomposition and measured costs.

\section{Related work}
Graph networks learn interactions\citep{battaglia2016}; GNS combines geometric particle graphs with message passing and noise-augmented training\citep{sanchez2020}. We build on released particle GNS software\citep{kumar2023}. Neural SPH incorporates physical structure into fluid prediction\citep{toshev2024}. Our intervention changes communication while retaining each trained simulator for a policy comparison.

MeshGraphNets predicts a sizing field for remeshing\citep{pfaff2021}, while LAMP learns refinement and coarsening under an error--computation tradeoff\citep{wu2023}. We retain the particle set and add local pairs. NRI and dNRI infer relational structure\citep{kipf2018,graber2020}; curvature-based rewiring addresses communication bottlenecks\citep{topping2022}; AdaMeshNet changes connectivity within message-passing layers\citep{seo2025}. These approaches motivate adaptation without equating prediction difficulty with the value of an additional edge. Faithful regression\citep{stirn2023} addresses the mean--variance optimization problem studied by \citet{seitzer2022}; we adopt its mean-gradient separation as a control.

\section{Limitations and conclusion}
The evidence comes from three-seed exploratory comparisons and includes continuations from inspected parents. Different data, horizons and backends do not isolate material effects; Sand's native-CUDA lineage does not establish optimization equivalence. Radius thresholds and top-budget selection can be discontinuous, and extra edges can amplify error. Neither the conditional theory nor lower MSE establishes stable physical simulation. Conservation, peak memory and optimized runtime remain unmeasured.

Constructive graphs make additional messages a controllable intervention. Exposure can teach a simulator to use those messages, but residual ranking does not consistently identify where they help. The contrast between Goop/WaterDrop exposure gains and Sand's adverse cached-risk feedback makes that distinction consequential. Adaptive graph methods should be judged by the benefit of their interventions and the trajectories they create, with complete policy comparisons alongside their computational cost.

\clearpage
\typeout{REVISION-MAIN-PAGES: \the\numexpr\value{page}-1\relax}
% After flushing the main text, the next page must be at most page 9.
% Fail compilation rather than silently exceed the eight-page main-text cap.
\ifnum\value{page}>9
  \PackageError{research-revision}{Main text exceeds the eight-page limit}{Shorten the main text before submission.}
\fi
\section*{AI Use Statement}
AI tools assisted with code, experiment execution, analysis, and manuscript editing.

\begin{thebibliography}{20}
\bibitem[Battaglia et~al.(2016)]{battaglia2016} P.~Battaglia et~al. Interaction networks for learning about objects, relations and physics. \emph{NeurIPS}, 2016.
\bibitem[Sanchez-Gonzalez et~al.(2020)]{sanchez2020} A.~Sanchez-Gonzalez et~al. Learning to simulate complex physics with graph networks. \emph{ICML}, 2020.
\bibitem[Kumar and Vantassel(2023)]{kumar2023} K.~Kumar and J.~Vantassel. GNS: A generalizable graph neural network-based simulator for particulate and fluid modeling. \emph{JOSS}, 8(88):5025, 2023. doi:10.21105/joss.05025.
\bibitem[Pfaff et~al.(2021)]{pfaff2021} T.~Pfaff et~al. Learning mesh-based simulation with graph networks. \emph{ICLR}, 2021.
\bibitem[Wu et~al.(2023)]{wu2023} Tailin Wu et~al. Learning controllable adaptive simulation for multi-resolution physics. \emph{ICLR}, 2023.
\bibitem[Kipf et~al.(2018)]{kipf2018} T.~Kipf et~al. Neural relational inference for interacting systems. \emph{ICML}, 2018.
\bibitem[Graber and Schwing(2020)]{graber2020} C.~Graber and A.~Schwing. Dynamic neural relational inference. \emph{CVPR}, 2020.
\bibitem[Topping et~al.(2022)]{topping2022} J.~Topping et~al. Understanding over-squashing and bottlenecks on graphs via curvature. \emph{ICLR}, 2022.
\bibitem[Seo et~al.(2025)]{seo2025} S.~Seo et~al. Adaptive graph rewiring to mitigate over-squashing in mesh-based GNNs for fluid dynamics simulations. arXiv:2511.12709, 2025.
\bibitem[Seitzer et~al.(2022)]{seitzer2022} M.~Seitzer, A.~Tavakoli, D.~Antic, and G.~Martius. On the pitfalls of heteroscedastic uncertainty estimation with probabilistic neural networks. \emph{ICLR}, 2022.
\bibitem[Stirn et~al.(2023)]{stirn2023} A.~Stirn et~al. Faithful heteroscedastic regression with neural networks. \emph{AISTATS}, PMLR 206:5593--5613, 2023.
\bibitem[Toshev et~al.(2024)]{toshev2024} A.~Toshev et~al. Neural SPH: Improved neural modeling of Lagrangian fluid dynamics. \emph{ICML}, PMLR 235:48428--48452, 2024.
\end{thebibliography}

% Official AISTATS 2027 checklist questions retained verbatim.
\section*{Checklist}

\begin{enumerate}

  \item For all models and algorithms presented, check if you include:
  \begin{enumerate}
    \item A clear description of the mathematical setting, assumptions, algorithm, and/or model. [Yes] The method, residual-risk interpretation, graph conventions, and assumptions of the conditional allocation bound are specified.
    \item An analysis of the properties and complexity (time, space, sample size) of any algorithm. [No] Graph cardinality, selection complexity and measured costs are provided; a complete complexity analysis for every implementation is not.
    \item (Optional) Anonymized source code, with specification of all dependencies, including external libraries. [Yes] Anonymized source code and dependency specifications accompany the paper.
  \end{enumerate}

  \item For any theoretical claim, check if you include:
  \begin{enumerate}
    \item Statements of the full set of assumptions of all theoretical results. [Yes] The stated propositions explicitly condition on their moment, feasibility, and additive-surrogate assumptions; no unconditional stability claim is made.
    \item Complete proofs of all theoretical results. [Yes] Proofs and derivations are included in the supplementary material.
    \item Clear explanations of any assumptions. [Yes] The text explains the limits of isotropic residual risk, lag assumptions, and the distinction between additive surrogate gain and nonlinear simulator benefit.
  \end{enumerate}

  \item For all figures and tables that present empirical results, check if you include:
  \begin{enumerate}
    \item The code, data, and instructions needed to reproduce the main experimental results (either in the supplemental material or as a URL). [No] The supplement regenerates figures and tables from supplied measurements and includes training and evaluation code. Complete datasets, full-model weights and every data-preparation path are not included.
    \item All the training details (e.g., data splits, hyperparameters, how they were chosen). [No] The appendix and protocols specify splits, model architecture and update schedules; the paper does not enumerate every optimizer and implementation setting.
    \item A clear definition of the specific measure or statistics and error bars (e.g., with respect to the random seed after running experiments multiple times). [Yes] Captions define position error, paired seed effects, sample standard deviations and undefined full-horizon outcomes.
    \item A description of the computing infrastructure used. (e.g., type of GPUs, internal cluster, or cloud provider). [Yes] Appendix~\ref{sec:implementation-details} specifies devices, software and graph backends; Appendix~\ref{sec:physical-cost} defines the timing scope.
  \end{enumerate}

  \item If you are using existing assets (e.g., code, data, models) or curating/releasing new assets, check if you include:
  \begin{enumerate}
    \item Citations of the creator If your work uses existing assets. [Yes] The original simulator software, public simulation data, and adopted regression and adaptive-simulation methods are cited.
    \item The license information of the assets, if applicable. [No] Software notices are included; dataset-specific redistribution terms are not fully documented.
    \item New assets either in the supplemental material or as a URL, if applicable. [Yes] Source code, experiment configurations and numerical results accompany the paper. Full-model weights are not included.
    \item Information about consent from data providers/curators. [Not Applicable] The experiments use existing publicly distributed synthetic particle-simulation data; no new human data were collected.
    \item Discussion of sensible content if applicable, e.g., personally identifiable information or offensive content. [Not Applicable] The studied inputs are synthetic particle states, not personal or offensive content.
  \end{enumerate}

  \item If you used crowdsourcing or conducted research with human subjects, check if you include:
  \begin{enumerate}
    \item The full text of instructions given to participants and screenshots. [Not Applicable] No crowdsourcing or human-subject experiment was conducted.
    \item Descriptions of potential participant risks, with links to Institutional Review Board (IRB) approvals if applicable. [Not Applicable] No crowdsourcing or human-subject experiment was conducted.
    \item The estimated hourly wage paid to participants and the total amount spent on participant compensation. [Not Applicable] No participant compensation was involved.
  \end{enumerate}

\end{enumerate}

\clearpage
% CURATED_APPENDIX_INSERT

\end{document}
